\section{Morfismos equidimensionais}
\label{section:IV.13}

\oldpage[IV-3]{187}

Esta seção é dedicada ao estudo da variação da \emph{dimensão} das fibras
de um morfismo localmente de tipo finito $f:X\to Y$ (que já apareceu a
propósito da ``fórmula das dimensões'' em~\sref{IV.5.5} e~\sref{IV.5.6}).
Demonstramos primeiro (\sref{IV.13.1.3}) que a função
$x\mapsto\dim_x f^{-1}(f(x))$ é sempre \emph{semicontínua superiormente}
em $X$ (teorema de semicontinuidade de Chevalley).
Estudamos em seguida, mais particularmente, os morfismos chamados
``equidimensionais'', para os quais essa função é \emph{localmente constante}.
\oldpage[IV-3]{188}
Infelizmente, a noção de morfismo equidimensional não é estável por mudança
de base; por isso, em muitas questões, é mais conveniente trabalhar com a
noção de morfismo universalmente aberto, cujo estudo é o objeto das
\textsection\textsection~14 e~15.

\subsection{O teorema de semicontinuidade de Chevalley}
\label{subsection:IV.13.1}

\begin{lemma}[13.1.1]
\label{IV.13.1.1}
Sejam $Y$ um pré-esquema irredutível localmente noetheriano,
$X$ um pré-esquema irredutível e $f:X\to Y$ um morfismo dominante de tipo finito.
Seja $\xi$ (resp.\ $\eta$) o ponto genérico de $X$ (resp.\ de $Y$),
e seja $e=\dim(f^{-1}(\eta))$.
Então, para todo $x\in X$, todas as componentes irredutíveis de
$f^{-1}(f(x))$ têm dimensão $\geq e$.
\end{lemma}

\begin{proof}
A afirmação é imediata quando, para todo $y\in f(X)$, $\sh{O}_y$ é um
anel \emph{universalmente catenário}: de fato, se $x$ é o ponto genérico
de uma componente irredutível $Z$ de $f^{-1}(y)$, de~\sref{IV.5.6.5},
combinado com~\sref{IV.5.2.1}, resulta que
$e+\dim(\sh{O}_y)=\dim(Z)+\dim(\sh{O}_x)$; mas, por~\sref[0]{0.16.3.9},
temos $\dim(\sh{O}_x)\leq\dim(\sh{O}_y)$, o que dá a conclusão nesse caso.

Vamos reduzir o caso geral a esse caso particular.
A questão é evidentemente local sobre $Y$ e, levando em conta
\sref{IV.4.1.1.3}, também é local sobre $X$; podemos, portanto, restringir-nos
ao caso em que $Y=\Spec(A)$ e $X=\Spec(B)$ são afins e irredutíveis,
sendo $B$ uma $A$-álgebra de tipo finito.
Além disso, por~\sref{IV.1.5.4}, podemos supor que $X$ e $Y$ são reduzidos,
de modo que $A$ e $B$ são domínios integrais.
Como $f$ é dominante, $A$ é então um \emph{subanel} de $B$.
Consideremos $A$ como limite indutivo de suas sub-$\bb{Z}$-álgebras
de tipo finito; de~\sref{IV.8.9.1} resulta então que existem uma tal
subálgebra $A_0$ e uma $A_0$-álgebra de tipo finito $B_0$ tais que
$B=B_0\otimes_{A_0}A$.
Ponhamos $Y_0=\Spec(A_0)$, $X_0=\Spec(B_0)$, e seja $f_0:X_0\to Y_0$
o morfismo correspondente ao homomorfismo $A_0\to B_0$, de modo que
$f=f_0\times 1_Y$.
Não é evidente \emph{a priori} que o pré-esquema $X_0$ seja integral,
mas veremos que podemos reduzir-nos a esse caso.
Seja $\eta_0$ o ponto genérico de $Y_0$; assim, se $g$ denota o morfismo
$Y\to Y_0$, temos $g(\eta)=\eta_0$.
Pela transitividade das fibras~\sref[I]{I.3.6.4}, temos
$f^{-1}(\eta)=f_0^{-1}(\eta_0)\otimes_{k(\eta_0)}k(\eta)$ e, como
$f^{-1}(\eta)$ é irredutível por hipótese, o mesmo vale para
$f_0^{-1}(\eta_0)$~\sref{IV.4.4.1}.
Nossa afirmação resultará então do lema seguinte:

\begin{lemma}[13.1.2]
\label{IV.13.1.2}
Sejam $Y_0$, $Y$ dois pré-esquemas integrais de pontos genéricos
$\eta_0$, $\eta$, seja $g:Y\to Y_0$ um morfismo dominante,
e seja $f_0:X_0\to Y_0$ um morfismo dominante tal que
$f_0^{-1}(\eta_0)$ seja irredutível.
Seja $X_0'$ a única componente irredutível de $X_0$ que encontra
$f_0^{-1}(\eta_0)$~\hyperref[0.2.1.8]{($\mathbf{0}_{\mathrm{I}}$,~2.1.8)},
e denotemos também por $X_0'$ o subpré-esquema fechado reduzido de $X_0$
cujo espaço subjacente é $X_0'$.
Suponhamos que o pré-esquema $X=X_0\times_{Y_0}Y$ seja integral;
então $X$ é isomorfo a $X_0'\times_{Y_0}Y$.
\end{lemma}

De fato, se $j_0:X_0'\hookrightarrow X_0$ é a injeção canônica,
que é uma imersão fechada,
$j=j_0\times 1_Y:X_0'\times_{Y_0}Y\to X_0\times_{Y_0}Y=X$
é uma imersão fechada.
Por outro lado, $X_0'$ contém a fibra $f_0^{-1}(\eta_0)$,
de modo que $X_0'\times_{Y_0}Y$ contém $f^{-1}(\eta)$;
observemos também que $f^{-1}(\eta)$ não é vazia~\sref[I]{I.3.4.7},
logo contém o ponto genérico de $X$.
Consequentemente, a imagem de $j$ é necessariamente todo $X$.
Mas, como $X$ é integral, o único subpré-esquema fechado de $X$
cujo espaço subjacente é $X$ é o próprio $X$;
portanto $j$ é um isomorfismo.

Estabelecido esse lema, podemos, portanto, supor que $X_0$ seja integral;
para todo $y_0\in Y_0$, $\sh{O}_{y_0}$ é uma $\bb{Z}$-álgebra
essencialmente de tipo finito, logo um anel universalmente
\oldpage[IV-3]{189}
catenário~\sref{IV.5.6.4}.
Consequentemente, para todo $y_0\in f_0(X_0)$, todas as componentes
irredutíveis de $f_0^{-1}(y_0)$\footnote{A fonte impressa traz aqui
$f_0^{-1}(y)$; o ponto quantificado e o argumento exigem $f_0^{-1}(y_0)$.}
têm dimensão $\geq e$, pois sabemos que
$\dim(f_0^{-1}(\eta_0))=e$ pela transitividade das fibras~\sref{IV.4.1.4}.
Para todo $y\in f(X)$, temos então $y_0=g(y)\in f_0(X_0)$ e,
pela transitividade das fibras e por~\sref{IV.4.2.7},
isso conclui a demonstração.
\end{proof}

\begin{theorem}[13.1.3]
\label{IV.13.1.3}
\textnormal{(Chevalley).} ---
Seja $f:X\to Y$ um morfismo localmente de tipo finito.
Para todo inteiro $n$, o conjunto $\mathrm{F}_n(X)$ dos $x\in X$
tais que $\dim_x(f^{-1}(f(x)))\geq n$ é fechado;
em outras palavras, a função $x\mapsto\dim_x(f^{-1}(f(x)))$
é semicontínua superiormente em $X$.
\end{theorem}

\begin{proof}
I) Suponhamos primeiro que $f$ seja localmente de apresentação finita.
A questão é evidentemente local sobre $X$ e sobre $Y$; podemos, portanto,
supor que $Y=\Spec(A)$ e $X=\Spec(B)$ são afins,
sendo $B$ uma $A$-álgebra de apresentação finita.
Sabemos então, por~\sref{IV.8.9.1}, que existem um subanel noetheriano $A_0$
de $A$ e uma $A_0$-álgebra de tipo finito $B_0$ tais que, pondo
$Y_0=\Spec(A_0)$ e $X_0=\Spec(B_0)$, temos
\[
X=X_0\times_{Y_0}Y,\qquad f=f_0\times 1_Y,
\]
onde $f_0:X_0\to Y_0$ é o morfismo correspondente ao homomorfismo $A_0\to B_0$.
Seja $g:Y\to Y_0$ o morfismo correspondente à injeção canônica $A_0\to A$,
seja $y$ um ponto de $Y$, e ponhamos $y_0=g(y)$;
sabemos que $f^{-1}(y)=f_0^{-1}(y_0)\otimes_{k(y_0)}k(y)$.
De~\sref{IV.4.2.7} resulta que, se $p$ é a projeção canônica
$f^{-1}(y)\to f_0^{-1}(y_0)$, as componentes irredutíveis de $f^{-1}(y)$
são as componentes irredutíveis dos espaços $p^{-1}(Z_0)$,
onde $Z_0$ percorre o conjunto das componentes irredutíveis de $f_0^{-1}(y_0)$;
cada componente irredutível de $p^{-1}(Z_0)$ domina $Z_0$
e tem dimensão $\dim(Z_0)$.
Levando em conta~\sref[0]{0.14.1.5}, vemos, portanto, que,
se $g':X\to X_0$ é a projeção canônica, $x$ um ponto de $X$ e $x_0=g'(x)$,
temos $\dim_x(f^{-1}(f(x)))=\dim_{x_0}(f_0^{-1}(f_0(x_0)))$,
de onde $\mathrm{F}_n(X)={g'}^{-1}(\mathrm{F}_n(X_0))$.
Reduzimo-nos, assim, a demonstrar o teorema quando $Y$ é \emph{noetheriano},
o que suporemos daqui em diante.

Podemos evidentemente supor que $X$ e $Y$ são reduzidos~\sref{IV.1.5.4}.
Considerando o conjunto das partes fechadas $Y'$ de $Y$ tais que o teorema
seja verdadeiro para o subpré-esquema fechado de $Y$ cujo espaço subjacente
é $Y'$ e para $X'=f^{-1}(Y')$, podemos raciocinar por indução
noetheriana~\hyperref[0.2.2.2]{($\mathbf{0}_{\mathrm{I}}$,~2.2.2)}
e supor que o teorema seja verdadeiro para toda parte fechada $Y'\neq Y$ de $Y$.
Se $X_i$ ($1\leq i\leq m$) são os subpré-esquemas fechados reduzidos de $X$
cujos espaços subjacentes são as componentes irredutíveis de $X$,
temos $\mathrm{F}_n(X)=\bigcup\mathrm{F}_n(X_i)$,
em virtude de~\sref[0]{0.14.1.5}.
Podemos, portanto, restringir-nos a demonstrar o teorema para cada $X_i$;
em outras palavras, podemos supor que $X$ seja \emph{irredutível}.
Se $Z$ é o subpré-esquema fechado de $Y$ cujo espaço subjacente é
$\overline{f(X)}$, $f$ se fatora como $X\xrightarrow{g}Z\xrightarrow{j}Y$,
onde $j$ é a injeção canônica~\sref[I]{I.5.2.2}
e $g$ é de tipo finito~\sref{IV.1.5.4}.
Basta, portanto, demonstrar o teorema para $Z$ e $g$.
Pela hipótese de indução, reduzimo-nos assim ao caso $Z=Y$,
isto é, ao caso em que $Y$ é irredutível e $f$ é dominante.
Seja então $\eta$ o ponto genérico de $Y$, e ponhamos $e=\dim(f^{-1}(\eta))$;
de~\sref{IV.13.1.1} resulta que, para $n\leq e$, temos $\mathrm{F}_n(X)=X$,
de modo que podemos restringir-nos ao caso em que $n>e$.
Mas então, por~\sref{IV.9.5.6}, existe uma vizinhança aberta $U$ de $\eta$
em $Y$ tal que $\mathrm{F}_n(X)\subset f^{-1}(Y-U)$.
Como $Y-U\neq Y$, a hipótese de indução implica que $\mathrm{F}_n(X)$ é fechado.

II) Passemos agora ao caso geral, escrevendo ainda $S=\Spec(A)$
para o alvo e supondo que $X=\Spec(B)$ seja afim,
com $B$ uma $A$-álgebra de tipo finito, portanto da forma
\oldpage[IV-3]{190}
$B=A[T_1,\ldots,T_n]/\mathfrak{J}$.
Seja $(\mathfrak{J}_\lambda)$ a família dos ideais finitamente gerados de
$A[T_1,\ldots,T_n]$ contidos em $\mathfrak{J}$,
de modo que $\mathfrak{J}$ é a união filtrante dos $\mathfrak{J}_\lambda$.
Se considerarmos $X$ e os
$X_\lambda=\Spec(A[T_1,\ldots,T_n]/\mathfrak{J}_\lambda)$
como subpré-esquemas fechados de $Z=\Spec(A[T_1,\ldots,T_n])$,
temos, para os espaços subjacentes, $X=\bigcap_\lambda X_\lambda$.
Se $f_\lambda:X_\lambda\to S$ é o morfismo estrutural,
deduzimos que $f^{-1}(s)=\bigcap_\lambda f_\lambda^{-1}(s)$
para todo $s\in S$.
Como os conjuntos $f_\lambda^{-1}(s)$ são fechados no espaço noetheriano $Z_s$,
existe um índice $\lambda$ (que depende de $s$)
tal que $f^{-1}(s)=f_\lambda^{-1}(s)$.
Para todo $x\in X$, ponhamos $d(x)=\dim_x(f^{-1}(f(x)))$ e
$d_\lambda(x)=\dim_x(f_\lambda^{-1}(f_\lambda(x)))$.
O que precede mostra que $d(x)=\inf_\lambda d_\lambda(x)$.
Como as funções $d_\lambda$ são semicontínuas superiormente
pela primeira parte da demonstração, o mesmo vale para $d$.
\end{proof}

\begin{corollary}[13.1.4]
\label{IV.13.1.4}
Sob as hipóteses de~\sref{IV.13.1.3}, o conjunto dos $x\in X$
tais que $x$ seja isolado em $f^{-1}(f(x))$ é aberto em $X$.
\end{corollary}

\begin{proof}
De fato, é o complementar de $\mathrm{F}_1(X)$~\sref[0]{0.14.1.10}.
\end{proof}

Observemos que recuperamos assim, sob hipóteses mais gerais,
a consequência~\sref[III]{III.4.4.10} do “Teorema principal” de Zariski.

\begin{corollary}[13.1.5]
\label{IV.13.1.5}
Sob as hipóteses de~\sref{IV.13.1.3}, suponhamos ainda que $f$
seja um morfismo fechado.
Então, para todo inteiro $n$, o conjunto dos $y\in Y$ tais que
$\dim(f^{-1}(y))\geq n$ é fechado; em outras palavras, a aplicação
$y\mapsto\dim(f^{-1}(y))$ é semicontínua superiormente.
Em particular, se $y$ é uma especialização de $y'$,
temos $\dim(f^{-1}(y))\geq\dim(f^{-1}(y'))$.
\end{corollary}

\begin{proof}
De fato, dizer que $\dim(f^{-1}(y))\geq n$ significa que
$y\in f(\mathrm{F}_n(X))$~\sref[0]{0.14.1.6}.
\end{proof}

\begin{corollary}[13.1.6]
\label{IV.13.1.6}
Sejam $X$, $Y$ dois pré-esquemas irredutíveis,
e seja $f:X\to Y$ um morfismo dominante localmente de tipo finito.
Seja $\xi$ (resp.\ $\eta$) o ponto genérico de $X$ (resp.\ de $Y$),
e seja $e=\dim(f^{-1}(\eta))$.
Então, para todo $x\in X$, temos $\dim_x(f^{-1}(f(x)))\geq e$.
\end{corollary}

\begin{proof}
De fato, o conjunto dos $x$ tais que $\dim_x(f^{-1}(f(x)))<e$
é aberto por~\sref{IV.13.1.3} e, como não pode conter $\xi$, é vazio.
\end{proof}

Observemos, por fim, o seguinte resultado mais simples:

\begin{proposition}[13.1.7]
\label{IV.13.1.7}
Seja $Y$ um pré-esquema quase-compacto,
e seja $f:X\to Y$ um morfismo de tipo finito.
Existe um inteiro $n$ tal que, para todo $y\in Y$,
temos $\dim(f^{-1}(y))\leq n$.
\end{proposition}

\begin{proof}
Como existe uma cobertura finita de $Y$ por abertos afins $(V_i)$
tal que cada $f^{-1}(V_i)$ é uma união finita de abertos afins,
reduzimo-nos imediatamente ao caso em que $Y=\Spec(A)$ e $X=\Spec(B)$
são afins, sendo $B$ uma $A$-álgebra de tipo finito.
Se $B$ admite um sistema de $n$ geradores, então, para todo $y\in Y$,
$B\otimes_A k(y)$ é uma $k(y)$-álgebra que admite $n$ geradores;
portanto $\dim(f^{-1}(y))\leq n$, em virtude de~\sref{IV.4.1.1}.
\end{proof}

\subsection{Morfismos equidimensionais: caso dos morfismos dominantes de pré-esquemas irredutíveis}
\label{subsection:IV.13.2}

\begin{env}[13.2.1]
\label{IV.13.2.1}
Sejam $Y$ um pré-esquema irredutível, $X$ um pré-esquema irredutível
e $f:X\to Y$ um morfismo \emph{dominante} localmente de tipo finito;
seja $\eta$ o ponto genérico de $Y$.
Sabemos, por~\sref{IV.13.1.6}, que, para todo $x\in X$,
temos $\dim_x(f^{-1}(f(x)))\geq\dim(f^{-1}(\eta))$.
\end{env}

\oldpage[IV-3]{191}
\begin{definition}[13.2.2]
\label{IV.13.2.2}
Sob as hipóteses de~\sref{IV.13.2.1}, dizemos que $f$ é
\emph{equidimensional no ponto} $x$ (ou que $X$ é \emph{equidimensional sobre}
$Y$ \emph{no ponto} $x$) se
\[
  \dim_x(f^{-1}(f(x)))=\dim(f^{-1}(\eta)).
\]
Dizemos que $f$ é \emph{equidimensional} (ou que $X$ é
\emph{equidimensional sobre} $Y$) se $f$ é equidimensional em todo ponto $x\in X$.
\end{definition}

Do teorema de Chevalley~\sref{IV.13.1.3} resulta que o conjunto dos
$x\in X$ nos quais $f$ é equidimensional é \emph{aberto e não vazio}.
Além disso, se $f$ é equidimensional no ponto $x$, todas as componentes
irredutíveis de $f^{-1}(f(x))$ que contêm $x$ têm a \emph{mesma dimensão},
pois cada uma tem dimensão $\geq\dim(f^{-1}(\eta))$~\sref{IV.13.1.6}
e $\leq\dim_x(f^{-1}(f(x)))$, em virtude de~\sref[0]{0.14.1.5}.

\begin{proposition}[13.2.3]
\label{IV.13.2.3}
Sejam $Y$ um pré-esquema irredutível localmente noetheriano,
$X$ um pré-esquema irredutível e $f:X\to Y$ um morfismo dominante
localmente de tipo finito.
Sejam $\eta$ o ponto genérico de $Y$, $x$ um ponto de $X$ e $y=f(x)$,
e suponhamos que se tenha
\[
\label{IV.13.2.3.1}
\tag{13.2.3.1}
  \dim(\sh{O}_x)=\dim(\sh{O}_y)+\dim(\sh{O}_x\otimes_{\sh{O}_y}k(y)).
\]
Então $f$ é equidimensional no ponto $x$.
A recíproca é verdadeira se houver igualdade em~\sref{IV.5.6.5.2},
em particular se $\sh{O}_y$ for universalmente catenário.
\end{proposition}

\begin{proof}
Isso resulta imediatamente de~\sref{IV.5.6.5.2}
e da desigualdade $\delta(x)\geq 0$~\sref{IV.13.1.1}.
\end{proof}

\begin{env}[13.2.4]
\label{IV.13.2.4}
Mais geralmente, sejam agora $Y$ um pré-esquema localmente noetheriano,
$f:X\to Y$ um morfismo de tipo finito, $x$ um ponto de $X$
e $X'$ uma parte fechada irredutível de $X$ que contém $x$;
ponhamos $y=f(x)$ e $Y'=\overline{f(X')}$,
e seja $\eta'$ o ponto genérico de $Y'$.
Denotemos também por $X'$ e $Y'$ os subpré-esquemas fechados reduzidos
de $X$ e $Y$ cujos espaços subjacentes são, respectivamente, $X'$ e $Y'$.
A restrição $X'\to Y$ de $f$ se fatora como
$X'\xrightarrow{f'}Y'\xrightarrow{j}Y$,
onde $j$ é a injeção canônica~\sref[I]{I.5.2.2}
e $f'$ é de tipo finito~\sref{IV.1.5.4}.
Ponhamos, para abreviar,
\[
\label{IV.13.2.4.1}
\tag{13.2.4.1}
  A=\sh{O}_{Y,y},\quad B=\sh{O}_{X,x},\quad A'=\sh{O}_{Y',y},\quad B'=\sh{O}_{X',x}.
\]
A fórmula~\sref{IV.5.6.5.2}, aplicada ao morfismo dominante $f'$
e aos pré-esquemas irredutíveis $X'$, $Y'$, dá
\[
\label{IV.13.2.4.2}
\tag{13.2.4.2}
  \dim(B')\leq\dim(A')+\dim(B'\otimes_{A'}k(y))-(\dim_x(f'^{-1}(y))-\dim(f'^{-1}(\eta')))
\]
Por outro lado, o anel local $A'$ (resp.\ $B'$, $B'\otimes_{A'}k(y)$)
é um quociente de $A$ (resp.\ de $B$, $B\otimes_A k(y)$);
portanto, por~\sref[0]{0.16.1.2.1}, temos
\[
\tag{13.2.4.3}
\label{IV.13.2.4.3}
  \dim(A')\leq\dim(A),\quad \dim(B')\leq\dim(B),\quad \dim(B'\otimes_{A'}k(y))\leq\dim(B\otimes_A k(y)).
\]
Deduzimos, portanto, em primeiro lugar, de~\sref{IV.13.2.4.3}
e de~\sref[0]{0.16.3.9},
\[
\label{IV.13.2.4.4}
\tag{13.2.4.4}
  \dim(B')\leq\dim(A')+\dim(B'\otimes_{A'}k(y))\leq\dim(A)+\dim(B\otimes_A k(y)).
\]
Além disso, em virtude de~\sref[0]{0.16.3.9},
temos também as desigualdades
\[
\label{IV.13.2.4.5}
\tag{13.2.4.5}
  \dim(B')\leq\dim(B)\leq\dim(A)+\dim(B\otimes_A k(y)).
\]
\end{env}

\oldpage[IV-3]{192}

Comparando essas desigualdades, obtemos:

\begin{lemma}[13.2.5]
\label{IV.13.2.5}
Com as notações de~\sref{IV.13.2.4}, as condições seguintes são equivalentes:
\begin{enumerate}[label=\emph{\alph*)}]
  \item $\dim(B')=\dim(A)+\dim(B\otimes_A k(y))$.
  \item Valem simultaneamente as relações seguintes:
  \begin{enumerate}[label=(\roman*)]
    \item $\dim(A')=\dim(A)$.
    \item $\dim(B'\otimes_{A'}k(y))=\dim(B\otimes_A k(y))$, em outras palavras,
    \[
      \dim_x(X'\cap f^{-1}(y))=\dim_x(f^{-1}(y)).
    \]
    \item $\dim_x(f'^{-1}(y))=\dim(f'^{-1}(\eta'))$,
    isto é, $f':X'\to Y'$ é equidimensional~\sref{IV.13.2.2} no ponto $x$.
    \item Temos a igualdade
    \[
      \dim(B')=\dim(A')+\dim(B'\otimes_{A'}k(y))-(\dim_x(f'^{-1}(y))-\dim(f'^{-1}(\eta')))
    \]
  \end{enumerate}
  (relação que sempre se verifica quando $A=\sh{O}_{Y,y}$ é um anel
  universalmente catenário, em virtude de~\sref{IV.5.6.5}).
  \item Valem simultaneamente as relações seguintes:
  \begin{enumerate}[label=(\roman*)]
    \item $\dim(B')=\dim(B)$.
    \item $\dim(B)=\dim(A)+\dim(B\otimes_A k(y))$.
  \end{enumerate}
\end{enumerate}
\end{lemma}

\begin{env}[13.2.6]
\label{IV.13.2.6}
Recordemos agora que as componentes irredutíveis $X_i$ de $X$ que contêm $x$
são em número finito e que, por~\sref{IV.5.1.2.1} e~\sref[0]{0.14.2.1.1},
\[
\label{IV.13.2.6.1}
\tag{13.2.6.1}
  \dim(\sh{O}_{X,x})=\sup_i\dim(\sh{O}_{X_i,x}).
\]
A equivalência das condições $b)$ e $c)$ de~\sref{IV.13.2.5} implica,
portanto, levando em conta~\sref[0]{0.16.3.9} e~\sref[0]{0.14.2.1}:
\end{env}

\begin{proposition}[13.2.7]
\label{IV.13.2.7}
Sejam $Y$ um pré-esquema localmente noetheriano,
$f:X\to Y$ um morfismo de tipo finito, $x$ um ponto de $X$, $f(x)=y$; temos
\[
\tag{13.2.7.1}
\label{IV.13.2.7.1}
  \dim(\sh{O}_{X,x})\leq\dim(\sh{O}_{Y,y})+\dim(\sh{O}_{X,x}\otimes_{\sh{O}_{Y,y}}k(y)).
\]
Para que os dois membros de~\sref{IV.13.2.7.1} sejam iguais, é necessário e
suficiente que exista um subconjunto fechado irredutível $X'$ de $X$ que
contenha $x$ e satisfaça simultaneamente as condições seguintes:
\begin{enumerate}[label=(\roman*)]
  \item Se $Y'=\overline{f(X')}$, temos
  $\dim(\sh{O}_{Y',y})=\dim(\sh{O}_{Y,y})$.
  \item $\dim_x(X'\cap f^{-1}(y))=\dim_x(f^{-1}(y))$
  (em outras palavras, $X'$ contém uma das componentes irredutíveis de
  $f^{-1}(y)$ de dimensão máxima entre as que contêm $x$).
  Isso equivale a dizer que
  $\dim(\sh{O}_{X',x}\otimes_{\sh{O}_{Y',y}}k(y))=\dim(\sh{O}_{X,x}\otimes_{\sh{O}_{Y,y}}k(y))$.
  \item $X'$ é equidimensional sobre $Y'$ no ponto $x$; em outras palavras,
  temos
  \[
    \dim_x(X'\cap f^{-1}(y))=\dim(X'\cap f^{-1}(\eta')),
  \]
  onde $\eta'$ é o ponto genérico de $Y'$ (e, portanto, todas as componentes
  irredutíveis de $X'\cap f^{-1}(y)$ que contêm $x$ têm a mesma
  dimensão~\sref{IV.13.2.2}).
\oldpage[IV-3]{193}
  \item Temos a igualdade
  \[
    \dim(\sh{O}_{X',x})=\dim(\sh{O}_{Y',y})+\dim(\sh{O}_{X',x}\otimes_{\sh{O}_{Y',y}}k(y))
  \]
  (condição sempre implicada por (iii) quando $\sh{O}_{Y,y}$ é um anel
  universalmente catenário).
\end{enumerate}
Além disso, $X'$ é então uma componente irredutível de $X$ e $Y'$ uma
componente irredutível de $Y$.
\end{proposition}

No enunciado, os anéis locais $\sh{O}_{X',x}$ e $\sh{O}_{Y',y}$ são relativos
aos subpré-esquemas fechados reduzidos de $X$, $Y$ cujos espaços subjacentes
são, respectivamente, $X'$ e $Y'$.

Além disso, isso também demonstra:

\begin{corollary}[13.2.8]
\label{IV.13.2.8}
Se os dois membros de~\sref{IV.13.2.7.1} são iguais, os subconjuntos fechados
irredutíveis $X'$ de $X$ que contêm $x$ e satisfazem as condições (i) a (iv)
de~\sref{IV.13.2.7} são exatamente aqueles para os quais temos
$\dim(\sh{O}_{X',x})=\dim(\sh{O}_{X,x})$ (o que implica necessariamente que
$X'$ é uma componente irredutível de $X$).
\end{corollary}

A Proposição~\sref{IV.13.2.7} implica:

\begin{corollary}[13.2.9]
\label{IV.13.2.9}
Sejam $Y$ um pré-esquema localmente noetheriano, $f:X\to Y$ um morfismo de
tipo finito, $x$ um ponto de $X$, $f(x)=y$.
Suponhamos que $\sh{O}_{Y,y}$ seja um anel universalmente catenário.
Então as condições seguintes são equivalentes:
\begin{enumerate}[label=\emph{\alph*)}]
  \item Os dois membros de~\sref{IV.13.2.7.1} são iguais.
  \item Existe uma componente irredutível $Z$ de $f^{-1}(y)$ que contém $x$,
  de dimensão $\dim_x(f^{-1}(y))$, tal que, para todo $x'$ em uma vizinhança
  de $x$ em $Z$, temos
  \[
  \label{IV.13.2.9.1}
\tag{13.2.9.1}
    \dim(\sh{O}_{X,x'})=\dim(\sh{O}_{Y,y})+\dim(\sh{O}_{X,x'}\otimes_{\sh{O}_{Y,y}}k(y)).
  \]
  \item Existe uma componente irredutível $Z$ de $f^{-1}(y)$ que contém $x$,
  de dimensão $\dim_x(f^{-1}(y))$, tal que, para o ponto genérico $z$ de $Z$,
  temos
  \[
  \label{IV.13.2.9.2}
\tag{13.2.9.2}
    \dim(\sh{O}_{X,z})=\dim(\sh{O}_{Y,y}).
  \]
\end{enumerate}
\end{corollary}

Mostremos que $a)$ implica $b)$.
Seja $\dim_x(f^{-1}(y))=n$; em virtude de~\sref{IV.13.2.7}, existe uma
componente irredutível $X'$ de $X$ que satisfaz as condições (i) a (iv)
de~\sref{IV.13.2.7}; seja $Z$ uma componente irredutível de dimensão $n$ de
$f^{-1}(y)$, que contém $x$ e está contida em $X'$.
Como $f^{-1}(y)$ é localmente noetheriano, existe uma vizinhança aberta $U$
de $x$ em $Z$ tal que $U$ não encontra nenhuma componente irredutível de
$f^{-1}(y)$ além das que contêm $x$; logo, por~\sref{IV.4.1.1.3},
$\dim_{x'}(f^{-1}(y))=n$ para todo $x'\in U$.
É então claro que as condições (i) a (iii) de~\sref{IV.13.2.7} são
satisfeitas quando substituímos $x$ por um ponto qualquer $x'\in U$, e o
mesmo vale para a condição (iv), pois $\sh{O}_{Y,y}$ é universalmente
catenário; daí a conclusão por~\sref{IV.13.2.7}.
A condição $b)$ implica trivialmente $c)$ em virtude de~\sref{IV.5.1.2}.
Finalmente, se $c)$ é satisfeita e se $X''$ é uma componente irredutível de
$X$ que contém $Z$ e tal que as condições (i) a (iv) de~\sref{IV.13.2.7}
são satisfeitas quando substituímos $X'$ por $X''$ e $x$ por $z$, é claro
que essas condições também são verificadas para $X''$ e $x$, pois
$\sh{O}_{Y,y}$ é universalmente catenário; portanto, $c)$ implica $a)$.

\begin{proposition}[13.2.10]
\label{IV.13.2.10}
Sejam $Y$ um pré-esquema irredutível localmente noetheriano, $\eta$ seu ponto
genérico, $f:X\to Y$ um morfismo de tipo finito, $y$ um ponto de $f(X)$.
Sejam $Z_i$ as
\oldpage[IV-3]{194}
componentes irredutíveis de $f^{-1}(y)$, $z_i$ o ponto genérico de $Z_i$, e
consideremos as condições seguintes:
\begin{enumerate}[label=\emph{\alph*)}]
  \item Para todo $x\in f^{-1}(y)$, temos a relação
  \[
  \tag{13.2.10.1}
\label{IV.13.2.10.1}
    \dim(\sh{O}_{X,x})=\dim(\sh{O}_{Y,y})+\dim(\sh{O}_{X,x}\otimes_{\sh{O}_{Y,y}}k(y)).
  \]
  \item Para todo $i$, temos
  \[
  \label{IV.13.2.10.2}
\tag{13.2.10.2}
    \dim(\sh{O}_{X,z_i})=\dim(\sh{O}_{Y,y}).
  \]
  \item Para todo $i$, existe uma componente irredutível $X_i$ de $X$ que
  contém $Z_i$ e tal que $\dim(X_i\cap f^{-1}(\eta))=\dim(Z_i)$
  (em outras palavras, o \emph{subpré-esquema fechado reduzido} $X_i$ de $X$
  é equidimensional sobre $Y$ no ponto $z_i$).
\end{enumerate}
Então $a)$ implica $b)$ e $b)$ implica $c)$; além disso, se $\sh{O}_{Y,y}$
é universalmente catenário, as três condições $a)$, $b)$, $c)$ são equivalentes.
\end{proposition}

Como o anel $\sh{O}_{X,z_i}\otimes_{\sh{O}_{Y,y}}k(y)$ é de dimensão $0$,
por~\sref{IV.5.1.2}, $a)$ implica evidentemente $b)$; $b)$ implica $c)$ em
virtude de~\sref{IV.13.2.7}, aplicada ao ponto $z_i$.
Reciprocamente, suponhamos que $\sh{O}_{Y,y}$ seja universalmente catenário;
como a condição $c)$ implica que $f(X_i)$ é denso em $Y$, segue de $c)$ que
as condições (i) a (iv) de~\sref{IV.13.2.7} são satisfeitas quando
substituímos $X'$ por $X_i$ e $x$ por $z_i$; logo, $c)$ implica $b)$;
finalmente, $b)$ implica $a)$ em virtude de~\sref{IV.13.2.9}.

\begin{corollary}[13.2.11]
\label{IV.13.2.11}
Com a notação de~\sref{IV.13.2.10}, suponhamos que $\sh{O}_{Y,y}$ seja
universalmente catenário.
Para todo $y'\in Y$, seja $E(y')$ o conjunto das dimensões das componentes
irredutíveis de $f^{-1}(y')$, e ponhamos
$d(y')=\dim(f^{-1}(y'))=\sup(E(y'))$.
Então, se as condições equivalentes $a)$, $b)$, $c)$ de~\sref{IV.13.2.10}
são satisfeitas, temos $E(y)\subset E(\eta)$, de onde $d(y)\leq d(\eta)$.
\end{corollary}

De fato, com a notação de~\sref{IV.13.2.10}, $X_i\cap f^{-1}(\eta)$ não é
vazio e é, portanto, uma componente irredutível de $f^{-1}(\eta)$,
por~\sref[0]{0.2.1.8}.

\begin{remarks}[13.2.12]
\label{IV.13.2.12}
\begin{enumerate}[label=(\roman*)]
  \item Recordemos, por~\sref{IV.6.1.2}, que a relação~\sref{IV.13.2.10.1}
  sempre se verifica quando o morfismo $f$ é \emph{plano} no ponto $x$.
  \item Com a notação de~\sref{IV.13.2.10}, suponhamos que $\sh{O}_{Y,y}$
  seja universalmente catenário e, além disso, que o morfismo $f$ seja
  \emph{próprio}; então, se as condições equivalentes $a)$, $b)$, $c)$
  de~\sref{IV.13.2.10} são satisfeitas, temos até mesmo $d(y)=d(\eta)$,
  pois~\sref{IV.13.1.5} implica $d(y)\geq d(\eta)$.
  \item O morfismo de~\sref{IV.12.2.3},~(ii) é \emph{próprio} e
  \emph{plano}, e todos os anéis locais de $Y$ são universalmente catenários;
  além disso, as duas componentes irredutíveis $X_1$, $X_2$ de $X$ são
  \emph{equidimensionais} sobre $Y$ em todo ponto; mas $E(y)$ tem
  \emph{dois} elementos para todo $y\neq y_0$, enquanto $E(y_0)$ se reduz
  a um \emph{único} elemento; portanto, $E(y)$ \emph{não é constante} em $Y$.
\end{enumerate}
\end{remarks}

\subsection{Morfismos equidimensionais: caso geral}
\label{subsection:IV.13.3}

\begin{proposition}[13.3.1]
\label{IV.13.3.1}
Sejam $Y$ um pré-esquema, $f:X\to Y$ um morfismo localmente de tipo finito,
$x$ um ponto de $X$, $y=f(x)$.
Denotemos por $Y_\alpha$ as componentes irredutíveis de $Y$ que contêm $y$.
Então as condições seguintes são equivalentes:
\begin{enumerate}[label=\emph{\alph*)}]
\oldpage[IV-3]{195}
  \item Existem um inteiro $e$ e uma vizinhança aberta $U$ de $x$ tais que
  a imagem por $f$ de toda componente irredutível de $U$ é densa em algum
  $Y_\alpha$ e que, para todo $x'\in U$, o espaço
  $U\cap f^{-1}(f(x'))$ é equidimensional e de dimensão $e$.
  \item[\emph{a$'$)}] Existem um inteiro $e$ e uma vizinhança aberta $U$
  de $x$ tais que a imagem por $f$ de toda componente irredutível de $U$
  é densa em algum $Y_\alpha$ e que, se denotamos por $y_\alpha$ o ponto
  genérico de $Y_\alpha$, toda componente irredutível dos espaços
  $U\cap f^{-1}(y)$, $U\cap f^{-1}(y_\alpha)$ é de dimensão $e$.
  \item[\emph{a$''$)}] Existem um inteiro $e$ e uma vizinhança aberta $U$
  de $x$ tais que, para cada uma das componentes irredutíveis $U_\lambda$
  de $U$, $f(U_\lambda)$ é denso em algum $Y_\alpha$ e que, para todo
  $x'\in U_\lambda$, as componentes irredutíveis de
  $U_\lambda\cap f^{-1}(f(x'))$ são todas de dimensão $e$.
  \item Existem um inteiro $e$, uma vizinhança aberta $U$ de $x$ e um
  $\mathrm Y$-morfismo quase-finito
  $g:U\to Y\otimes_{\bb{Z}}\bb{Z}[T_1,\ldots,T_e]$
  (pré-esquema que também denotaremos por $Y[T_1,\ldots,T_e]$, para abreviar)
  tais que a imagem por $g$ de toda componente irredutível de $U$ é densa
  em uma componente irredutível de $Y[T_1,\ldots,T_e]$.
\end{enumerate}
\end{proposition}

\begin{proof}
É imediato que $a'')$ implica $a)$, pois, para todo $x\in U$, cada componente
irredutível de $U\cap f^{-1}(f(x))$ é uma componente irredutível de um dos
$U_\lambda\cap f^{-1}(f(x))$; daí a conclusão por~\sref[0]{0.14.1.4}.
A condição $a)$ implica trivialmente $a')$.
Mostremos que $a')$ implica $a'')$; podemos nos restringir ao caso em que
$f$ é de tipo finito e $X$ e $Y$ são reduzidos, por~\sref{IV.1.5.4}.
Seja $U_\lambda$ uma componente irredutível de $U$, e suponhamos que
$f(U_\lambda)$ seja denso em $Y_\alpha$; então a restrição de $f$ a
$U_\lambda$ se fatora como
$U_\lambda\xrightarrow{f_\alpha}Y_\alpha\longrightarrow Y$, onde $f_\alpha$
é de tipo finito e dominante, por~\sref[I]{I.5.2.2}.
Seja $x_\lambda$ o ponto genérico de $U_\lambda$; em virtude
de~\sref[0]{0.2.1.8}, $U_\lambda\cap f_\alpha^{-1}(y_\alpha)$ é a única
componente irredutível de $U\cap f^{-1}(y_\alpha)$ que contém $x_\lambda$
e, por hipótese, é de dimensão $e$, igual às dimensões de todas as
componentes irredutíveis de $U_\lambda\cap f^{-1}(y)$, em virtude da
hipótese e de~\sref{IV.13.1.1}.
Mas, pelo teorema de Chevalley~\sref{IV.13.1.3} e por~\sref{IV.13.1.1},
o conjunto $V_\lambda$ dos $x'\in U_\lambda$ tais que
$\dim_{x'}(f_\lambda^{-1}(f_\lambda(x')))=e$ é aberto e contém $x$;
basta tomar a união dos $V_\lambda$ para obter um conjunto aberto
que satisfaz as condições de $a'')$.

Provemos agora que $a)$ implica $b)$; podemos nos restringir ao caso em que
$Y=\Spec(A)$ e $X=\Spec(B)$ são afins e $U=X$.
Provemos primeiro o lema seguinte:

\begin{lemma}[13.3.1.1]
\label{IV.13.3.1.1}
Sejam $Y=\Spec(A)$, $X=\Spec(B)$ dois esquemas afins, $f:X\to Y$ um morfismo
de tipo finito, $y$ um ponto de $f(X)$, e ponhamos $e=\dim(f^{-1}(y))$.
Então existe um $\mathrm Y$-morfismo $g:X\to Y[T_1,\ldots,T_e]$ tal que
(se pomos $X_y=X\times_Y\Spec(k(y))$), o morfismo
$g_y:X_y\to\Spec(k(y)[T_1,\ldots,T_e])$ é finito.
Além disso, para todo morfismo $g$ desse tipo, $g_y$ é necessariamente
sobrejetivo e existe uma vizinhança aberta $U$ de $f^{-1}(y)$ em $X$
tal que $g|U$ é quase-finito.
\end{lemma}

\begin{proof}
Seja $\mathfrak{p}=\mathfrak{j}_y$; o anel
$B\otimes_A k(\mathfrak{p})$ é uma $k(\mathfrak{p})$-álgebra de tipo finito;
portanto, pelo lema de normalização (Bourbaki, \emph{Alg.\ comm.}, cap.~V,
\textsection~3, n$^{\text{o}}$~1, teorema~1), existem em
$B\otimes_A k(\mathfrak{p})$ uma sequência finita $(t_i)_{1\leq i\leq r}$
de elementos algebricamente independentes sobre $k(\mathfrak{p})$ e tais
que, se pomos $C'=k(\mathfrak{p})[t_1,\ldots,t_r]$,
$B\otimes_A k(\mathfrak{p})$ é uma $C'$-álgebra \emph{finita}; temos, portanto,
$\dim(B\otimes_A k(\mathfrak{p}))=\dim(C')$, por~\sref[0]{0.16.1.5},
e, como $\dim(C')=r$, por~\sref{IV.5.2.1}, temos $r=e$.
Como $B\otimes_A k(\mathfrak{p})=(B/\mathfrak{p}B)\otimes_{A/\mathfrak{p}}k(\mathfrak{p})$,
podemos, multiplicando os $t_i$ por um elemento $\neq0$
\oldpage[IV-3]{196}
conveniente de $A/\mathfrak{p}$, supor que cada $t_i$ é a imagem canônica
em $B\otimes_A k(\mathfrak{p})$ de um elemento $s_i\in B$.
Seja então $u:A[T_1,\ldots,T_e]\to B$ o homomorfismo tal que $u(T_i)=s_i$
para todo $i$, e seja $g:X\to Y[T_1,\ldots,T_e]$ o morfismo correspondente.
É claro que, pela escolha dos $s_i$, $g_y$ é um morfismo finito.
Para todo morfismo $g:X\to Y[T_1,\ldots,T_e]$ tal que $g_y$ seja finito,
de~\sref{IV.5.4.2} e~\sref{IV.4.1.2.1} resulta que $g_y$ é necessariamente
sobrejetivo.
Por outro lado, em virtude de~\sref{IV.13.1.4}, o conjunto $U$ dos $x\in X$
que são isolados em sua fibra $g^{-1}(g(x))$ é aberto em $X$ e contém
$f^{-1}(y)$; por definição, a restrição $g|U$ é um morfismo quase-finito.
\end{proof}

Estabelecido esse lema, para provar que $a)$ implica $b)$, resta mostrar
que, se $x_\lambda$ é um ponto maximal de $U$, então
$g(x_\lambda)=z_\lambda$ é um ponto maximal de $Z=Y[T_1,\ldots,T_e]$.
Em virtude da hipótese $a)$, podemos (restringindo $U$ se necessário) supor
que $f(x_\lambda)$ é um dos pontos genéricos $y_\alpha$ das componentes
irredutíveis de $Y$ que contêm $y$.
Se $p:Z\to Y$ é o morfismo estrutural, temos $p(z_\lambda)=y_\alpha$;
portanto, $g$ induz um $k(y_\alpha)$-morfismo quase-finito
\[
h:U\cap f^{-1}(y_\alpha)\longrightarrow p^{-1}(y_\alpha).
\]
Ora, $p^{-1}(y_\alpha)=\Spec(k(y_\alpha)[T_1,\ldots,T_e])$ é integral e de
dimensão $e$.
Se $z_\lambda$ não fosse um ponto maximal de $Z$, tampouco seria um ponto
maximal de $p^{-1}(y_\alpha)$, por~\sref[0]{0.2.1.8}, e seu fecho
$Z'_\lambda$ em $p^{-1}(y_\alpha)$ teria, portanto, dimensão $<e$,
por~\sref{IV.4.1.2.1}.
Mas, como $h$ é quase-finito, de~\sref{IV.4.1.2} e da hipótese $a)$ resulta
que $\dim(Z'_\lambda)\geq e$ (a restrição de $h$ a
$\overline{\{x_\lambda\}}$ se fatorando como
\[
\overline{\{x_\lambda\}}\longrightarrow Z'_\lambda\longrightarrow p^{-1}(y_\alpha)
\]
em virtude de~\sref[I]{I.5.2.2}).
Chegamos assim a uma contradição, o que demonstra que $a)$ implica $b)$.

Provemos finalmente que $b)$ implica $a)$.
Observemos que o morfismo estrutural $p:Z=Y[T_1,\ldots,T_e]\to Y$ é fielmente
plano; portanto, por~\sref{IV.2.3.4}, as imagens por $p$ dos pontos maximais
de $Z$ são os pontos maximais de $Y$.
Isso já demonstra que os $f(x_\lambda)$ são pontos genéricos dos $Y_\alpha$.
Além disso, se $f(x_\lambda)=y_\alpha$, o $k(y_\alpha)$-morfismo
$h:U\cap f^{-1}(y_\alpha)\to p^{-1}(y_\alpha)$ deduzido de $g$ é dominante
e quase-finito por hipótese; temos então, por~\sref{IV.4.1.2},
$\dim(U_\lambda\cap f^{-1}(y_\alpha))=e$ (sendo $U_\lambda$ a componente
irredutível de $U$ de ponto genérico $x_\lambda$).
Analogamente, para todo $x'\in U_\lambda$, o morfismo
$U_\lambda\cap f^{-1}(f(x'))\to p^{-1}(f(x'))$ deduzido de $g$ é um
$k(f(x'))$-morfismo quase-finito; logo, por~\sref{IV.4.1.2}, temos
$\dim(f^{-1}(f(x'))\cap U_\lambda)\leq e$.
Mas, por~\sref{IV.13.1.6}, todas as componentes irredutíveis de
$U_\lambda\cap f^{-1}(f(x'))$ são de dimensão $\geq e$; vemos assim que
essas componentes são exatamente de dimensão $e$; portanto, $b)$ implica $a)$.
\end{proof}

\begin{definition}[13.3.2]
\label{IV.13.3.2}
Sejam $Y$ um pré-esquema, $f:X\to Y$ um morfismo localmente de tipo finito,
$x$ um ponto de $X$.
Dizemos que $f$ é \emph{equidimensional no ponto} $x$ (ou que $X$ é
\emph{equidimensional sobre} $Y$ \emph{no ponto} $x$) se as condições
equivalentes de~\sref{IV.13.3.1} são satisfeitas.
Dizemos que $f$ é \emph{equidimensional} (ou que $X$ é
\emph{equidimensional sobre} $Y$) se $f$ é equidimensional em todo
ponto $x\in X$.
\end{definition}

É claro que, para que $f$ seja equidimensional em um ponto $x$, é necessário
e suficiente que $f_{\mathrm{red}}$ o seja.
Além disso, as condições de~\sref{IV.13.3.1} mostram que o conjunto dos
pontos em que $f$ é equidimensional é \emph{aberto} em $X$.

Observemos que, quando $X$ e $Y$ são \emph{irredutíveis}, dizer que $f$ é
\oldpage[IV-3]{197}
equidimensional no ponto $x$ significa, em virtude de~\sref{IV.13.3.1},~$a'')$,
que $f$ é dominante e que
$\dim_x(f^{-1}(f(x)))=\dim(f^{-1}(\eta))$ (onde $\eta$ é o ponto genérico
de $Y$); a definição~\sref{IV.13.3.2} coincide, portanto, neste caso, com
a definição~\sref{IV.13.2.2}.

\begin{proposition}[13.3.3]
\label{IV.13.3.3}
Sejam $Y$ um pré-esquema localmente noetheriano,
$f:X\to Y$ um morfismo localmente de tipo finito, $x$ um ponto de $X$
e $X_j$ as componentes irredutíveis de $X$ (em número finito) que contêm $x$.
Para que $f$ seja equidimensional no ponto $x$, é necessário e suficiente
que, para todo $j$, $Y_j=\overline{f(X_j)}$ seja uma componente irredutível
de $Y$ e que, denotando por $X_j$ e $Y_j$ os subpré-esquemas fechados
reduzidos de $X$ e $Y$ de espaços subjacentes $X_j$ e $Y_j$, por
$f_j:X_j\to Y_j$ o morfismo deduzido de $f$, por~\sref[I]{I.5.2.2},
e por $y_j$ o ponto genérico de $Y_j$, $f_j$ seja equidimensional no ponto
$x$ e todos os números $\dim(f_j^{-1}(y_j))$ sejam iguais.
\end{proposition}

\begin{proof}
Isso resulta imediatamente de~\sref{IV.13.3.1},~$a')$).
\end{proof}

\begin{corollary}[13.3.4]
\label{IV.13.3.4}
Com a notação de~\sref{IV.13.3.3}, ponhamos $y=f(x)$.
Se $\sh{O}_{Y,y}$ é um anel universalmente catenário e $f$ é
equidimensional no ponto $x$, temos
\[
\tag{13.3.4.1}
\label{IV.13.3.4.1}
  \dim(\sh{O}_{X_j,x})=\dim(\sh{O}_{Y_j,y})+e-\deg.\mathrm{tr}_{k(y)}k(x)
\]
onde $e$ é o valor comum dos números $\dim(f_j^{-1}(y_j))$.
\end{corollary}

\begin{proof}
Como cada um dos $\sh{O}_{Y_j,y}$, quociente de $\sh{O}_{Y,y}$, é um anel
universalmente catenário, por~\sref{IV.5.6.1}, a igualdade é consequência
de~\sref{IV.5.6.5}.
\end{proof}

\begin{corollary}[13.3.5]
\label{IV.13.3.5}
Com a notação de~\sref{IV.13.3.4}, suponhamos que $f$ seja equidimensional
no ponto $x$ e que $\sh{O}_{Y,y}$ seja um anel universalmente catenário.
Se o anel $\sh{O}_{Y,y}$ é equidimensional, o mesmo vale para
$\sh{O}_{X,x}$; a recíproca é verdadeira se a imagem por $f$ da união dos
$X_j$ é densa em uma vizinhança de $y$.
\end{corollary}

\begin{proof}
De fato, dizer que $\sh{O}_{Y,y}$ é equidimensional significa que os números
$\dim(\sh{O}_{Y'_i,y})$ são iguais para todas as componentes irredutíveis
$Y'_i$ de $Y$ que contêm $y$, como resulta de~\sref{IV.5.1.1.5},
aplicada ao esquema local $\Spec(\sh{O}_{Y,y})$; como temos a
relação~\sref{IV.13.3.4.1}, segue pelo mesmo raciocínio que
$\sh{O}_{X,x}$ é então equidimensional.
Reciprocamente, se $\sh{O}_{X,x}$ é equidimensional, todos os números
$\dim(\sh{O}_{Y_j,y})$ são iguais\footnote{A impressão omite aqui o operador
de dimensão, embora fale em números; restaura-se esse operador, como exige
a igualdade (13.3.4.1).} por~\sref{IV.13.3.4.1}; deduzimos que $\sh{O}_{Y,y}$
é equidimensional se os $Y_j$ são todas as componentes irredutíveis de $Y$
que contêm $y$, o que resulta da hipótese suplementar.
\end{proof}

\begin{proposition}[13.3.6]
\label{IV.13.3.6}
Sejam $Y$ um pré-esquema localmente noetheriano,
$f:X\to Y$ um morfismo localmente de tipo finito, $x$ um ponto de $X$,
$y=f(x)$.
Suponhamos que o anel $\sh{O}_y$ seja equidimensional.
Então, se $\sh{O}_x$ é equidimensional e temos a igualdade
\[
\tag{13.3.6.1}
\label{IV.13.3.6.1}
  \dim(\sh{O}_x)=\dim(\sh{O}_y)+\dim(\sh{O}_x\otimes_{\sh{O}_y}k(y))
\]
(cf.~\sref{IV.13.2.7.1}), $f$ é equidimensional no ponto $x$, e a recíproca
é verdadeira se $\sh{O}_y$ é um anel universalmente catenário.
\end{proposition}

\begin{proof}
Mantendo a notação de~\sref{IV.13.3.3}, de~\sref{IV.13.2.8} e da hipótese
de que $\sh{O}_x$ é equidimensional resulta que cada $Y_j$ é uma componente
irredutível de $Y$ e que cada $f_j$ é equidimensional no ponto $x$; além disso,
por~\sref{IV.13.2.8}, temos (levando em conta~\sref{IV.5.6.5})
\[
  \dim(\sh{O}_{X_j,x})=\dim(\sh{O}_{Y_j,y})+\dim_x(X_j\cap f^{-1}(y))-\deg.\mathrm{tr}_{k(y)}k(x).
\]

\oldpage[IV-3]{198}

Ora, como $\sh{O}_y$ é suposto equidimensional, essa igualdade se escreve
\[
\tag{13.3.6.2}
\label{IV.13.3.6.2}
  \dim_x(X_j\cap f^{-1}(y))=\dim(\sh{O}_{X,x})-\dim(\sh{O}_{Y,y})+\deg.\mathrm{tr}_{k(y)}k(x).
\]
O primeiro membro de~\sref{IV.13.3.6.2} é, portanto, independente de $j$;
mas, como $f_j$ é equidimensional no ponto $x$, temos
$\dim_x(X_j\cap f^{-1}(y))=\dim(f_j^{-1}(y_j))$; logo, o critério
de~\sref{IV.13.3.3} mostra que $f$ é equidimensional no ponto $x$.

Reciprocamente, suponhamos que $\sh{O}_{Y,y}$ seja um anel universalmente
catenário e que $f$ seja equidimensional no ponto $x$; então,
por~\sref{IV.13.3.5}, $\sh{O}_{X,x}$ é equidimensional, e de~\sref{IV.13.3.4}
resulta que temos a relação
\[
  \dim(\sh{O}_{X,x})=\dim(\sh{O}_{Y,y})+e-\deg.\mathrm{tr}_{k(y)}k(x)
\]
onde $e$ é o valor comum dos números
$\dim(f_j^{-1}(y_j))=\dim_x(X_j\cap f^{-1}(y))$; mas, por definição,
$e$ também é igual a $\dim_x(f^{-1}(y))$; daí a relação~\sref{IV.13.3.6.1},
levando em conta~\sref{IV.5.6.5.2}.
\end{proof}

\begin{proposition}[13.3.7]
\label{IV.13.3.7}
Sejam $Y$ um pré-esquema, $f:X\to Y$ um morfismo localmente de tipo finito
e equidimensional, e $Z$ um subconjunto fechado de $X$.
Então a função
\[
\label{IV.13.3.7.1}
\tag{13.3.7.1}
  x\mapsto\operatorname{codim}_x(Z\cap f^{-1}(f(x)),f^{-1}(f(x)))
\]
é semicontínua inferiormente em $X$.
\end{proposition}

\begin{proof}
Como todas as componentes irredutíveis de $f^{-1}(f(x))$ que contêm $x$
têm a mesma dimensão por hipótese, por~\sref{IV.13.3.1}, temos, em virtude
de~\sref{IV.5.2.1},
\[
  \operatorname{codim}_x(Z\cap f^{-1}(f(x)),f^{-1}(f(x)))=\dim_x(f^{-1}(f(x)))-\dim_x(Z\cap f^{-1}(f(x))).
\]
Mas, por hipótese, o primeiro termo do segundo membro é uma função contínua
de $x$, por~\sref{IV.13.3.1}, e o segundo é uma função semicontínua
superiormente de $x$, por~\sref{IV.13.1.3}\footnote{A fonte impressa cita
aqui (13.3.3); a semicontinuidade superior invocada é o teorema de
Chevalley (13.1.3).}; daí a conclusão.
\end{proof}

\begin{proposition}[13.3.8]
\label{IV.13.3.8}
Sejam $Y$ um pré-esquema, $f:X\to Y$ um morfismo localmente de tipo finito,
$x$ um ponto de $X$.
Sejam $Y'$ um pré-esquema, $g:Y'\to Y$ um morfismo plano,
$X'=X\times_Y Y'$, $f'=f_{(Y')}:X'\to Y'$; se $f$ é equidimensional no
ponto $x$, então $f'$ é equidimensional em todo ponto $x'\in X'$ situado
sobre $x$.
\end{proposition}

\begin{proof}
Como a questão é local em $X$ e $Y$, podemos nos restringir ao caso em que
toda componente irredutível de $X$ (resp.~de $Y$) contém $x$ (resp.~$y$);
como a imagem por $f$ de toda componente irredutível de $X$ é então densa
em uma componente irredutível de $Y$, por~\sref{IV.13.3.1}, sabemos,
por~\sref{IV.2.3.5}, que a imagem por $f'$ de toda componente irredutível
de $X'$ é densa em uma componente irredutível de $Y'$.
Por outro lado, pela transitividade das fibras, por~\sref[I]{I.3.6.4},
de~\sref{IV.4.2.8} e~\sref{IV.13.3.1} resulta que, para todo $z'\in X'$,
o conjunto das dimensões das componentes irredutíveis de
${f'}^{-1}(f'(z'))$ é o mesmo que o conjunto das dimensões das componentes
irredutíveis de $f^{-1}(f(z))$, onde $z$ é a projeção de $z'$ sobre $X$;
daí a conclusão em virtude de~\sref{IV.13.3.1},~$a)$).
\end{proof}

\begin{remarks}[13.3.9]
\label{IV.13.3.9}
A propriedade de equidimensionalidade de um morfismo $f:X\to Y$ \emph{não}
é estável por mudança de base arbitrária $g:Y'\to Y$, nem mesmo quando $g$
é a injeção canônica de uma componente irredutível $Y'$ de $Y$ em $Y$
(sendo $Y'$
\oldpage[IV-3]{199}
considerado como subpré-esquema fechado reduzido de $Y$).
Por exemplo, seja $k$ um corpo, seja $A_0=k[S,T]$ o anel de polinômios em
duas indeterminadas e ponhamos $\mathfrak{a}=\mathfrak{p}_1\mathfrak{p}_2$,
onde $\mathfrak{p}_1$ e $\mathfrak{p}_2$ são os ideais primos $A_0S$ e
$A_0T$ de $A_0$; sejam $A=A_0/\mathfrak{a}$ e $Y=\Spec(A)$, que tem duas
componentes irredutíveis $Y_1=\Spec(A/\mathfrak{p}_1)$ e
$Y_2=\Spec(A/\mathfrak{p}_2)$; tomemos $X=Y_1$, com $f:X\to Y$ a injeção
canônica, que é evidentemente um morfismo equidimensional.
Por outro lado, tomemos $Y'=Y_2$, com $g:Y'\to Y$ a injeção canônica; então
$X'=X\times_Y Y'=\Spec\bigl((A/\mathfrak{p}_1)\otimes_A(A/\mathfrak{p}_2)\bigr)=\Spec\bigl(A/(\mathfrak{p}_1+\mathfrak{p}_2)\bigr)$,
e $\mathfrak{p}_1+\mathfrak{p}_2$ é um ideal maximal de $A$, de modo que
$X'$ se reduz a um ponto; o morfismo $f':X'\to Y'$ \emph{não é dominante},
pois a imagem por $f'$ do único ponto de $X'$ é um ponto fechado de $Y'$;
portanto, $f'$ não é equidimensional.

Também podemos dar um contraexemplo em que $X$ e $Y$ são \emph{integrais},
$f$ é \emph{finito} e \emph{birracional} (e \emph{a fortiori} equidimensional,
por~\sref{IV.13.3.1},~$b)$), e $g:Y'\to Y$ é \emph{finito} e \emph{dominante}.
Sejam $A$ e $\overline{A}$ os anéis locais definidos em~\sref[II]{II.7.5},
e tomemos $Y=\Spec(A)$ e $X=\Spec(\overline{A})$; por outro lado, com a
notação de~\sref[II]{II.7.5}, tomemos $Y'=\Spec(\overline{B})$; então
$X'=\Spec(\overline{A}\otimes_A\overline{B})=\Spec(\overline{B}\otimes_B\overline{B})$;
mas verificamos imediatamente que $\overline{B}\otimes_B\overline{B}$ é
o produto direto dos anéis $B/\mathfrak{p}'$, $B/\mathfrak{p}''$ e de dois
anéis isomorfos a $B/\mathfrak{n}$, cujos espectros se reduzem, portanto,
a um ponto; como as projeções desses pontos são pontos fechados de $Y'$,
vemos aqui novamente que $f'$ não leva uma componente irredutível de $X'$
a um subconjunto denso de uma componente irredutível de $Y'$;
portanto, $f'$ não é equidimensional.

Neste exemplo, o anel $A$ não é geometricamente unirramificado; veremos no
parágrafo seguinte {\hyperlink{ega4.c04.cor.14-4-6}{(14.4.6)}} que tais
fenômenos não podem ocorrer quando os pontos de $Y$ são
\emph{geometricamente unirramificados}.
A falta de estabilidade da noção de morfismo equidimensional restringe
muito seu interesse, em favor da noção de morfismo universalmente aberto,
que será estudada em detalhe no parágrafo seguinte.
\end{remarks}
